\documentclass[a4paper,11pt]{memoir}
%\documentclass[b5paper,10pt,twoside]{memoir}
% Ved a4 format skal nedenst�ende nok justeres en del
\setlrmarginsandblock{3.5cm}{2.5cm}{*}
\setulmarginsandblock{3cm}{*}{1.2}

\usepackage[retainorgcmds]{IEEEtrantools}
\usepackage{graphicx}
\usepackage{pstricks,pst-node,pst-text,pst-3d}
\usepackage{marginnote}
\usepackage{array,booktabs,longtable}
\usepackage{url}
\usepackage{color}
\definecolor{webblue}{rgb}{0.0,0.05,0.45}
\definecolor{MyDarkBlue}{rgb}{0,0.08,0.45}
\definecolor{webred}{rgb}{0.75,0,0}

\usepackage[T1]{fontenc}
\checkandfixthelayout
% ops�tning af pagestyles
\makepagestyle{book} % laver en ny tom pagestyle
\makeevenhead{book}{}{\small\sffamily\leftmark}{}
\makeoddhead{book}{}{\small\sffamily\rightmark}{}
\makeevenfoot{book}{\small\sffamily\thepage}{}{}
\makeoddfoot{book}{}{}{\small\sffamily\thepage}
\makeatletter

\makepsmarks{book}{%
\renewcommand\chaptermark[1]{%
\markboth{%
\ifnum \value{secnumdepth} > 1
\if@mainmatter % indenfor frontmatter er der intet kapitel nummer
\@chapapp\ \thechapter. \ % \@chapapp er lidt dum, se nedenfor
\fi
\fi
##1}{}}%
\renewcommand\tocmark{\markboth{\contentsname}{\contentsname}}%
\renewcommand\lofmark{\markboth{\listfigurename}{\listfigurename}}%
\renewcommand\lotmark{\markboth{\listtablename}{\listtablename}}%
\renewcommand\bibmark{\markboth{\bibname}{\bibname}}%
\renewcommand\indexmark{\markboth{\indexname}{\indexname}}%
\renewcommand\sectionmark[1]{\markright{##1}}%
\renewcommand\subsectionmark[1]{\markright{##1}}%
\renewcommand\subsubsectionmark[1]{\markright{##1}}%
}
\makeatother

\pagestyle{book}
% laver om p� plain stilen
\copypagestyle{plain}{book}
\makeevenhead{plain}{}{}{}
\makeoddhead{plain}{}{}{}

%**************** Headlines ****************************************************
% konfiguration af kapitel titel font samt fonte til sections
\renewcommand\chapnamefont{\Huge\bfseries\sffamily}
\renewcommand\chapnumfont{\chapnamefont}
\renewcommand\chaptitlefont{\color{webred}\Huge\usefont{OT1}{phv}{bc}{n}\selectfont\raggedright}
\setsubsecheadstyle{\color{MyDarkBlue}\large\bfseries\sffamily\raggedright}
\setsubsubsecheadstyle{\color{webred}\Huge\usefont{OT1}{phv}{bc}{n}\selectfont\raggedright}
% Rule under headline.. (Underline section headlines med de 4 linier herunder) 
%\newcommand{\ruledsec}[1]{%
%\Large\usefont{OT1}{phv}{b}{n}\selectfont\raggedright #1 %\color{webred}\rule[15pt]{\textwidth}{1.0pt}} 
%\setsecheadstyle{\ruledsec} %ud-kommenter linie herunder
\setsecheadstyle{\Large\usefont{OT1}{phv}{b}{n}\selectfont\raggedright}

\usepackage{titlesec}
\usepackage{titletoc}
% Afstande mellem overskrifter og tekst
\titleformat{\subsection}
{\color{MyDarkBlue}\usefont{OT1}{phv}{b}{n}\selectfont}{\thesubsection}{1em}{}
\titleformat{\subsubsection}
{\color{MyDarkBlue}\usefont{OT1}{phv}{b}{n}\selectfont}{\thesubsubsection}{1em}{}
\titlespacing{\subsection}{15pt}{15pt}{10pt}
\titlespacing{\subsubsection}{15pt}{15pt}{10pt}
\titlespacing{\section}{15pt}{35pt}{10pt}
%***********  section number in margin ****************************************
\makeatletter
\def\@seccntformat#1{\@ifundefined{#1@cntformat}%
{\csname the#1\endcsname\quad}% default
{\csname #1@cntformat\endcsname}% individual control
}
\def\section@cntformat{\color{webred}\protect\makebox[0pt][r]{\thesection.\quad}}
\makeatother


% s�g for at hvis en \section el.lign. flyttes til en ny side
% da skal siden f�r ikke str�kkes
\raggedbottomsectiontrue
% justering af afsnitsnummerering og ToC dybde
\setsecnumdepth{subsubsection} % til og med
\maxsecnumdepth{subsubsection} % underligt koncept
\settocdepth{subsection} %til og med
% nogle pakker man ofte anvender
\usepackage[ansinew]{inputenc} % eller ansinew hvis Windows
%*******************************************************************************
\usepackage[T1]{fontenc}
\usepackage{amsmath,amssymb}
\usepackage{mathtools}
\usepackage{graphicx}
\usepackage{fix-cm,fixltx2e}
\usepackage{soul}

\DeclareRobustCommand{\SetFourierSpace}{%
\fontdimen2\font=1.13\fontdimen2\font}
\sodef\an{}{0.13em}{0em}{0em} \sodef\ann{}{0.13em}{0.5em}{0em}
%******************************************************************************
\let\footruleskip\relax % for compatibility of memoir and fancyhdr
%% Use fancy chapter headers, with Jos Dingjan's modifications,
%% plus my own tweaks. This style is not part of teTeX,
%% so we are using a local (and renamed) copy.
\usepackage[Lenny]{fncychapleo}
\usepackage{fancyhdr}
%*******************************************************************************

\makeatletter
\newcommand\figcaption{\def\@captype{figure}\caption}
\newcommand\tabcaption{\def\@captype{table}\caption}
\makeatother

\newenvironment{narrow}[2]{%
\begin{list}{}{%
\setlength{\topsep}{0pt}%
\setlength{\leftmargin}{#1}%
\setlength{\rightmargin}{#2}%
\setlength{\listparindent}{\parindent}%
\setlength{\itemindent}{\parindent}%
\setlength{\parsep}{\parskip}}%
\item[]}{\end{list}}

\newlength{\marginwidth}
\setlength{\marginwidth}{2.0\oddsidemargin} %bredden af billeder ud i margin'en

% Picture Handling
\usepackage{subfig}
\usepackage[leftcaption]{sidecap}
\usepackage{varioref} %For smarte referencer med \vref istedet for \ref
\usepackage{calc}% auto udregn

\graphicspath{{figures/}}
\sidecaptionvpos{figure}{t} %sidecaption aligned med toppen af billede
\usepackage{caption}
\DeclareCaptionFont{red}{\color{webred}}
\captionsetup{singlelinecheck=false,font=footnotesize,labelfont={bf,red},format=hang}

\DeclareCaptionFormat{llap}{\llap{#1#2}#3\par}
\strictpagechecktrue
\makeatletter
\DeclareRobustCommand*{\bfseries}{%
  \not@math@alphabet\bfseries\mathbf
  \fontseries\bfdefault\selectfont
  \boldmath
}
\makeatother

% andre pakker og konfigurationer
\usepackage{makeidx}
\usepackage{mflogo} 
\makeindex % hvis man anvender s�dan en


\renewcommand*{\cftchaptername}{Chapter\space}
\renewcommand*{\cftfigurename}{Fig.\space}
\renewcommand{\contentsname}{Table of Contents}
\renewcommand*\abstractname{Summary}

%\includeonly{introduction} %Hvis man ikke vil kompilere det hele hver gang
\begin{document}

%\begin{document}
\begin{titlingpage} \vspace*{5mm} \centering
%\huge \textsc{\ann{Johnny Rieper}} \par
%\small \an{2007\,0114} \vspace{3mm} \par
%\fontsize{1.2cm}{1cm}\selectfont \MakeUppercase{\an{Simulations on}}
\vspace{3mm} \par

%\fontsize{1.2cm}{1cm}\selectfont \MakeUppercase{\an{Matching Techniques}}
\Huge \MakeUppercase{\an{Distributed Computation}} \vspace{3mm} \par
\Large \MakeUppercase{\an{in }} \vspace{5mm} \par
\LARGE \MakeUppercase{\an{a Production Environment}} \vspace{3mm} \par
\vspace{3mm} \par
%\Large {\an{Introducing Hidden Markov Models \&}}  \par
%\Large {\an{Dynamic Time Warping}} \vspace{5mm} \par

%\fontsize{1.0cm}{1cm}\selectfont \MakeUppercase{\an{her}}
%\vspace{5mm} \par
\large \sodef\ann{}{0.13em}{0.5em}{0em}
\small \textsc{\ann{by}} \par
\vspace{2mm} \par
\Large \textsc{\ann{Theis Jeppe Nickelsen}} \par
\small \an{201300249} \vspace{3mm} \par
\Large \textsc{\ann{Stefan Krause-Kj\ae r}} \par
\small \an{201300256} \vspace{3mm} \par
\Large \textsc{\ann{Thomas Thisgaard Steffensen}} \par
\small \an{201300250} \vspace{3mm} \par
\enlargethispage{15mm} \vspace{20mm} \par

%\begin{figure}[t]
%\raggedleft
%\includegraphics[trim=0 0 0 0,clip,keepaspectratio,width=5cm]{redlogo.eps} \par
%\end{figure}


%\begin{figure}[h]
%\centering
%\includegraphics[width=3.0cm]{segla1b} \par
%\end{figure}

\begin{figure}[b]
%\centering
\begin{pspicture}(5,5)
\rput[bl](0,0){\colorbox{webred}{\parbox{\textwidth}{\color{white} \qquad \qquad \qquad \qquad \qquad \qquad \large\textsc{\ann{\textbf{Aarhus University}}}}}}
\rput[bl](.6,-0.51){\pscirclebox*[framesep=-10pt]{\includegraphics[scale=1]{figures/segla1s}}}
%\psset{linewidth=0.1,arrows=->,arrowscale=2}%
%\psline(6,3)(4,4.6)\uput[-45](6,3){Ear}
%\psline(6,1.5)(2,1.5)\uput[0](6,1.5){Mouth}
\end{pspicture}
\end{figure}

\vspace{15mm} \vfill \par

\normalsize\textsc{\ann{Master's Thesis}}\par
\vspace{2mm}
\normalsize\textsc{\ann{Supervisor: Christian Fischer Pedersen}}\par
\vspace{2mm}
\normalsize \scshape \sodef\ann{}{0.13em}{0.6em}{0em}
{\ann{Thesis on Distributed Computing}} \par
\vspace{2mm}
\ann{Aarhus School of Engineering} \par
%\ann{University of Aarhus} \par
\vspace{2mm}
\an{September 2014}

%\colorbox{webred}{\color{webblue}text}
\vspace{0mm} \par


\newpage
\thispagestyle{empty} % fjerne evt. sidehoved og fod
% \small forstyrrer ikke da vi er inden i et environment
\small
\strut\vfill % pres alt ned i bunden af siden
\copyright\ Theis Jeppe Nickelsen, Stefan Krause-Kj\ae r, Thomas Thisgaard Steffensen 2014\\[\baselineskip]
This Thesis is printed with Computer Modern 10pt\\
Layout and typography by the author using \LaTeX\\
%All simulated graphs designed by author in \textsc{Matlab}\texttrademark\\
%Remaining figures designed by author using PSTricks.

\end{titlingpage}


%\endinput
%\end{document}

\frontmatter

\end{document}



